\section{Prelude to Development}

\begin{quotex}
All of Saint Bernard's life seems destined to show that in order to solve problems of an intellectual and even a political order there exist means quite other than those we have long since become accustomed to considering the only ones effective, no doubt because they are the only ones within reach of a purely human wisdom, \emph{which is not even a shadow of true wisdom}.
\flright{\textsc{Rene Guenon}, \emph{Saint Bernard}}
\end{quotex}


One of the aims of this discussion of Solovyov is to demonstrate how a particular tradition treats the general questions of Tradition as \textbf{Rene Guenon} has shown. We pointed out that the three basic forms of being, that is, as we actually experience our life,--feeling, thinking, and willing---are associated with the Platonic Triad of fundamental ideas: Beauty, Truth, and Goodness. These ideas manifest on three levels: material, formal, and absolute (above form). As we shall see, that yields nine different areas of manifestation. Note that the three levels correspond to the traditional castes: Brahman, Kshatriya, Vaishya, which reflect the orientation of a person. Since a person may be centered on the feeling, thinking, or willing function, Solovyov's system demonstrates nine personality types. For example, one may focus on political matters, another on economic activity, another on art or mystical states. Each one will probably think his own area of interest is the most important and will have impatience with those who believe otherwise. The metaphysican, however, is interested in knowledge of the Whole, the One, as only this knowledge is true Wisdom.

That Solovyov's approach is viable, is shown by a similar notion from Guenon:

\begin{quotex}
Where something truly traditional is concerned, everything must be present from the beginning and subsequent developments serve only to render it more explicit without the adjunction of new and external elements.
\end{quotex}


This is from Guenon's important essay, The Holy Grail. Obviously, this is equivalent to Solovyov's definition of development. So we shall be interested in the development of Traditional ideas. For one, Solovyov's system explains why there are three castes. He will also deal with the question of spiritual authority and temporal power as it has developed. Finally, the whole cycle must be understood.

In particular, the conclusion we must draw from this is that the Tradition espoused by Solovyov is a subsequent development of the Primordial Tradition. The process of its development from the Satya Yuga to the Kali Yuga must be explained. Since development must have a goal, which cannot be total chaos because it is impossible, the goal is the recovery of the \textbf{Primordial State}.

\paragraph{Misunderstandings}

In our discussion of the development of Tradition, it is important to clear up some misunderstandings about what Tradition is not.

\begin{description}

\item[Documentarian]

Tradition is not like a documentary. While it may be useful to document how symbols or rites are used across traditions, that is not the main point. Knowing, for example, some ancient Etruscan fertility rite is not ipso facto knowing Tradition.
\item[Antiquarian]

Tradition is not an antiquarian interest. For example, the recovery of medieval Lithuanian basket weaving techniques is not Tradition. Arts and crafts may be one of the areas we discuss, but it is far from the Whole.
\item[Partisan]

Tradition is certainly not partisanship. To evangelize for one tradition over another is really an exoteric activity. However, it lacks the dispassion and objectivity required for an understanding of Tradition. In normal circumstances, a man follows the traditions of his people. That is why Guenon acted as a Catholic in France, but relocated to Egypt to be a Muslim.
\end{description}


\paragraph{Understandings}


Guenon makes several points that will help in an analysis of Tradition.

\begin{description}

\item[Privileged Form]

Guenon writes:

\begin{quotex}
There is no traditional form that is privileged; the only distinction to be made is between forms that have disappeared and those still living.
\end{quotex}


There is a marked tendency today to claim the superiority of one form over another, usually on the basis that it is more suitable to the mentality of a particular ethnic group. However, this is a serious misunderstanding, as it is concerned with the psychological. Guenon writes:

\begin{quotex}
True mystical states elude the domain of psychology entirely.
\end{quotex}

\item[Esoteric vs Exoteric]

When we discuss traditions, we are generally concerned with its esoteric aspects. Thus any criticisms based on particular exoteric manifestations of teachings, are irrelevant to us, and only reveal the ignorance of the critic. Hence, we agree with Guenon when he writes:

\begin{quotex}
There can be no doubt of the existence of Christian esoterism in the Middle Ages.
\end{quotex}


That is objectively and infallibly true. Anyone who disagrees understands neither esoterism nor the Middle Ages, and is typically deluded by psychological and sentimental motivations.

Many are confused about this, so we dare to hope that these quotes from Guenon will dispel any future confusions:

\begin{quotex}
Esoterism can under no circumstances be represented by churches of sects of any kind ... those who persist in ascribing to sects what concerns esoterism or initiation are on the wrong track and can only go astray.
\end{quotex}


This is not said in opposition to exoterism, but rather not to confuse domains. For example, exoterically Christians claim that \textbf{Jesus} is the Avatar of the \textbf{Logos}. His claim to being God does not rest on the circumstances of his birth,--which would make him a demigod--- but rather on the claim that the Logos is God. As a contingent, historical event, there is no metaphysical proof of this birth. Now, Solovyov claims to prove that the Logos is God, and even that the Logos must incarnate. However, he cannot prove that the incarnation happened at a certain time and place to a particular individual. Please keep this distinction in mind.
\item[Continuity and Incorporation]

As we have repeated many times, as one tradition fades away, what is essential in it is incorporated into another living tradition. As one illustration, Guenon writes:

\begin{quotationx}
When a traditional form is on the verge of extinction, its last representatives may very well deliberately entrust to this collective memory i.e., folklore what would otherwise be irrevocably lost. This is the only way to save what can be saved; and, at the same time, the natural incomprehension of the masses is a sufficient guarantee that whatever possesses an esoteric character will not be despoiled in the process but will remain as a sort of witness to the past for those in later times who may be capable of understanding it.

It seems beyond doubt that \emph{the origins of the Grail legend must be linked to the transmission from Druidism to Christianity of traditional elements of an initiatic order}. Once this transmission had been effected in a regular manner, whatever the modalities of the transmission may have been, \emph{these elements thereby became an integral part of Christian esoterism}.
\end{quotationx}


There is no question here of ``borrowing'', or ``deceiving'', or even distorting. The Grail legend is preserved in its fullness, becoming an integral part of the new tradition. Guenon never anticipated that the masses would have access to all this material; unfortunately, the incomprehension remains, but now results in despoiling. Thus we see neo-Druids and neo-pagans reappear, and like the \textbf{Aghoris}, feed on the cadavers of extinct traditions, and ignore the fresh food of still living traditions. In our example, if they cannot understand the Legend of the Grail as expressed in Medieval literature, how can they possibly recover is alleged original meaning?
\end{description}


\paragraph{Conclusion}


Guenon lists two works, \textbf{Dante}'s Divine Comedy and the Romance of the Rose, that describe the esoteric process quite openly. With that in mind, serious men will choose to study those works assiduously rather than indulge in fantasies. Guenon also points to Saint Bernard as an antidote to the poison of modernity:

\begin{quotationx}
Among the great figures of the Middle Ages, there are few whose study is more suited for counteracting certain prejudices cherished by the modern mind that Saint Bernard.

Saint Bernard was always determined by the same intentions:

\begin{itemize}


\item To defend the right

\item To combat injustice

\item To maintain unity in the Christian world
\end{itemize}
\end{quotationx}


That is enough to get started.


\flrightit{Posted on 2011-12-05 by Cologero}

\begin{center}* * *\end{center}

\begin{footnotesize}
\begin{sffamily}
\texttt{Avery on 2011-12-06 at 02:34 said:}

I don't think the new right is in danger of falling apart. It's a movement made up of intelligent people who lived liberalism and woke up to realize, ``this is insane''. When you step out of the Matrix you find a bunch of people disagreeing with each other over how to fight it, but you can at least understand their intentions. People who want spiritual disunity are fighting a different battle entirely.

Yes, I understand their intentions, not always their tactics. Our only concern here is to counter some common misconceptions about Tradition. Admin

\hfill

\texttt{EXIT on 2011-12-06 at 10:36 said:}

Guenon concluded that a restoration of Christianity was impossible, given its modernist leanings and as so meany perennialists agree today so much of Christianity (i.e., its gnostic esoteric elements) have been lost and cannot be put back again. Ere it is the christians who are the aghoris feeding on a corpse of a dead tradition. Christianity is particularly difficult and undesirable for the very reason that the church is not transmitter of the esoteric tradition as it should be were there anyone of any understanding and worth in the church. For without the holy of holies it is precisely useless and of zero interest to traditionalists.

As for the issue of borrowing... how can it be denied given all the evidence that christians borrowed from other traditions? Didn't even Evola admit that the christians distorted the grail legend? I suggest reading Murdock's The Christ Conspiracy. Christianity grew out of a multiculturalized world with Greek and Roman occupations of Palestine and the Alexandrian school. Whether these elements were synthesized or not is an entirely separate issue from that of borrowing; but as we often see, these elements were not always perfectly fitted together precisely because they were borrowings.

\hfill

\texttt{logres on 2011-12-06 at 22:45 said:}

I don't understand your quarrel with Christianity -- most of the Middle Ages was a repudiation of precisely the ``multicultural and synthesizing'' element which you decry in the decaying Roman world. I understand the dismay at modern Church, however much I might disagree about precisely how to fix it.

\hfill

\texttt{Charlotte Cowell on 2011-12-07 at 03:55 said:}

Exit, I do understand that it's impossible to get something until you get it. On this subject specifically Rudolf Steiner's lecture about the return of Christ in the Etheric is helpful. For those with an interest in Shambhala this lecture also has some interesting information towards the bottom of the article (which is quite short):

\url{http://wn.rsarchive.org/Lectures/GA118/English/APC1961/19100306p02.html}


\hfill

\texttt{Perennial on 2011-12-08 at 03:35 said:}

Are we in for a another mythinformation fest, EXIT?

\hfill

\texttt{Cologero on 2022-12-05 at 09:01 said:}

Well, the ``New Right'' did fall apart. And most of the self-appointed leaders turned out to be opportunists or even decadent. A ``bunch of people disagreeing'' does not make a ``movement''. There is a fundamental logical difference between A and not A as Guenon points out in Man and his Becoming.


\hfill

\end{sffamily}
\end{footnotesize}

